\documentclass[11pt]{article}
\usepackage[utf8]{inputenc}
\usepackage[T1]{fontenc}
\usepackage[a4paper,margin=1in]{geometry}
\usepackage{lmodern}
\usepackage{microtype}
\usepackage{amsmath,amssymb,mathtools}
\usepackage{mathrsfs}
\usepackage{tikz-cd}
\usepackage{enumitem}
\usepackage[hidelinks]{hyperref}
\setlength{\parindent}{0pt}
\setlength{\parskip}{0.55em}
\allowdisplaybreaks
\newcommand{\Z}{\mathbb Z}
\newcommand{\Q}{\mathbb Q}
\newcommand{\Ahat}{\widehat A}
\newcommand{\Atil}{\widetilde A}
\newcommand{\Ctil}{\widetilde C}
\newcommand{\K}{\mathrm K}
\newcommand{\Gr}{\operatorname{G}}
\newcommand{\ch}{\operatorname{ch}}
\newcommand{\rank}{\operatorname{rank}}
\newcommand{\Tor}{\operatorname{Tor}}
\newcommand{\Coker}{\operatorname{Coker}}
\newcommand{\Ker}{\operatorname{Ker}}
\newcommand{\Supp}{\operatorname{Supp}}
\newcommand{\codim}{\operatorname{codim}}
\newcommand{\cO}{\mathcal O}
\newcommand{\cC}{\mathcal C}
\newcommand{\cF}{\mathscr F}
\newcommand{\cE}{\mathscr E}
\newcommand{\eps}{\varepsilon}
\newcommand{\ellc}{\ell}
\begin{document}

\section*{SGA 6, Appendix to Expos\'e 0: RRR}
\begin{center}
{\Large Chapter II, Sections 1--5}\par
\vspace{0.4em}
{\large Classes of coherent algebraic sheaves and Chern classes}\par
\vspace{0.4em}
Source scan pages 47--75, printed pages 40--68.
\end{center}

Thus, relative to proper morphisms, and ignoring multiplicative structures and gradings, $A(X)$ is a covariant functor in $X$. Let us record the classical formula
\begin{equation}\tag{2.3}
 f_*(x\cdot f^*(y))=f_*(x)\cdot y .
\end{equation}

We shall also use Chow's theory of Chern classes. It consists in attaching to every vector bundle $E$ on $X$ (still assumed quasi-projective and nonsingular) Chern classes
\[
 c^i(E)\in A^i(X)\qquad (i\geq 1),
\]
in such a way as to satisfy the following conditions. Put $c^0(E)=1$ and
\begin{equation}\tag{2.4}
 \Ctil(E)=\left[\rank E,\sum_{i\geq 0}c^i(E)\right]\in \Atil(X)
\end{equation}
(see Chapter I, \S 3 for the definition of $\Atil(X)$). The conditions to be verified by a theory of Chern classes are then expressed as follows:
\begin{equation}\tag{2.5}
 \Ctil(E)=\Ctil(E')+\Ctil(E'')
\end{equation}
if the vector bundle $E$ is an extension of $E'$ by $E''$. If one puts
\[
 C(E)=\sum_i c^i(E),
\]
this condition is expressed by $C(E)=C(E')C(E'')$. Formula (2.5) permits one to extend $\Ctil$ to an additive homomorphism from $K(\xi(X))$ (cf. Chapter I, \S 2, Example 2) to $\Atil(X)$:
\begin{equation}\tag{2.6}
 \Ctil:K(\xi(X))\longrightarrow \Atil(X),
\end{equation}
and one wants $\Ctil$ to be a $\lambda$-homomorphism of $\lambda$-rings. This last condition may also be written
\begin{equation}\tag{2.7}
 \Ctil(E\otimes F)=\Ctil(E)\Ctil(F),\qquad
 \Ctil(\lambda^iE)=\lambda^i\Ctil(E),
\end{equation}
for two vector bundles $E$ and $F$; of course $\Ctil$ is automatically compatible with augmentations. Let $D$ be a divisor on $X$ and let $L(D)$ be the vector bundle it defines. One wants
\begin{equation}\tag{2.8}
 c^1(L(D))=\ellc(D),\qquad c^i(L(D))=0\quad (i>1),
\end{equation}
where, for any cycle $Z$, $\ellc(Z)$ denotes its class in $A(X)$. Finally, one imposes that the $c^i$ be functorial, i.e. that if $f$ is a morphism from $X$ to $Y$ and $E$ is a vector bundle on $Y$, then
\begin{equation}\tag{2.9}
 c^i(f^!(E))=f^*(c^i(E)),
\end{equation}
where $f^!(E)$ is the inverse image of $E$ on $X$. One may also write (2.9) in the following form: the notion of inverse image of vector bundles defines a $\lambda$-homomorphism of augmented $\lambda$-rings
\begin{equation}\tag{2.10}
 f^!:K(Y)\longrightarrow K(X),\footnote{The notation $f^!$ used here is in conflict with that used in duality theory (cf., for example, R. Hartshorne, \emph{Residues and Duality}, Lecture Notes in Mathematics, no. 20, Springer, 1966). This is why in this seminar we replace it by the notation $f^*$.}
\end{equation}
where here one writes $K(X)$ instead of $K(\xi(X))$, and (2.9) means that there is commutativity in the diagram
\[
\begin{tikzcd}[column sep=large,row sep=large]
K(Y) \arrow[r,"f^!"] \arrow[d,"\Ctil_Y"'] & K(X) \arrow[d,"\Ctil_X"]\\
A(Y) \arrow[r,"f^*"'] & A(X).
\end{tikzcd}
\]
Here the space on which $\Ctil$ is considered has been put as a subscript.

\textbf{Theorem 2.2.} There exists a theory of Chern classes satisfying conditions (2.5), (2.7), (2.8), and (2.9).

We shall also admit the following proposition.

\textbf{Proposition 2.3.} Let $Y$ be a closed subset of $X$ and let $U$ be its complement. If $x$ is an element of $A(X)$ inducing $0$ on $U$, then there exists a representative of $x$ which is a cycle carried by $Y$.

\textbf{Corollary.} Let $p=\dim X-\dim Y$. Then, if $i<p$, the homomorphism
\[
 A^i(X)\longrightarrow A^i(U)
\]
is bijective. Every element of the kernel of $A^p(X)\to A^p(U)$ is a linear combination of components of dimension $n-p$ of $Y$; hence it is proportional to $\ellc(Y)$ if $Y$ is irreducible.

\subsection*{2. Definition of the Chern classes of coherent algebraic sheaves}

Let $\cC$ be an abelian category, and let $\cC_0$ be a subclass of $\cC$. Denote by $K(\cC_0)$ the quotient group of the free group generated by the isomorphism classes of elements of $\cC_0$ by the subgroup generated by the elements of the form
\[
(E)-(E')-(E''),
\]
where $E$ is an extension of $E'$ by $E''$ and $E,E',E''\in\cC_0$. We must suppose that the classes of elements of $\cC_0$, for the isomorphism relation, form a set. Denote by $\gamma_{\cC_0}(E)$ the class of an object $E\in\cC_0$ in $K(\cC_0)$, and omit the subscript $\cC_0$ when there is no danger of confusion. The group $K(\cC_0)$ is the solution of a universal problem relative to functions
\[
E\longmapsto f(E)
\]
defined on $\cC_0$, with values in an arbitrary abelian group $A$, which are additive in the sense that $f(E)=f(E')+f(E'')$ whenever $E$ is an extension of $E'$ by $E''$.

Let $F$ be an additive functor from $\cC$ to $\cC'$ and let $\cC'_0$ be a subclass of $\cC'$ satisfying the same condition as $\cC_0$, such that for every $A\in\cC_0$ the $R^qF(A)$ are in $\cC'_0$ and are zero for $q$ sufficiently large. We suppose that injective resolutions exist in $\cC$, so that the derived functors $R^qF$ of $F$ exist. Then $F$ defines a homomorphism from $K(\cC_0)$ to $K(\cC'_0)$, still denoted $F$, by the formula
\begin{equation}\tag{2.11}
 F(\gamma_{\cC_0}(E))=\sum_{i\geq 0}(-1)^i\gamma_{\cC'_0}(R^iF(E)).
\end{equation}
More generally, if one has a cohomological functor $(F^i)$ from $\cC$ to $\cC'$ such that for every $E\in\cC_0$ the $F^i(E)$ are in $\cC'_0$ and are zero except for a finite number of $i$, one obtains a homomorphism from $K(\cC_0)$ to $K(\cC'_0)$. These remarks extend to multifunctors; moreover, the classical spectral sequences give transitivity properties which we shall not make explicit.

In particular, if $F$ is an exact functor, formula (2.11) simplifies to
\begin{equation}\tag{2.11 bis}
 F(\gamma_{\cC_0}(E))=\gamma_{\cC'_0}(F(E)).
\end{equation}
This applies in particular if $\cC=\cC'$, if $F$ is the identity functor, and if $\cC_0\subset \cC'_0$; one has a canonical homomorphism from $K(\cC_0)$ to $K(\cC'_0)$.

\textbf{Theorem 2.3.} Let $\cC$ be an abelian category and let $\cC_0\subset \cC_1$ be two subclasses such that the isomorphism classes of objects of $\cC_1$ form a set. Suppose that the following conditions are satisfied:

\begin{enumerate}[label=(\roman*)]
\item For every exact sequence $0\to A'\to A\to A''\to 0$ in $\cC$ such that $A$ and $A''$ are in $\cC_0$ (respectively $\cC_1$), one has $A'\in\cC_0$ (respectively $A'\in\cC_1$).
\item Let $u:A\to B$ and $v:C\to B$, with $A,B,C\in\cC_1$ and $u$ surjective. Then there exist $L\in\cC_0$ and homomorphisms $u':L\to C$ and $v':L\to A$ such that $u'$ is surjective and $vu'=uv'$.
\item For every $A\in\cC_1$ there exists an integer $d=d(A)$ such that for every left resolution $L$ of $A$ by objects of $\cC_0$ one has $Z_d(L)\in\cC_0$, where $Z_d$ denotes the cycles in dimension $d$.
\end{enumerate}
Under these conditions, the canonical homomorphism from $K(\cC_0)$ to $K(\cC_1)$ is an isomorphism.

\emph{Principle of the proof.} Let $A\in\cC_1$. By (ii) and (iii) there exists a finite resolution $L$ of $A$ by objects of $\cC_0$; put
\[
 \ell(A)=\sum_{i\geq 0}(-1)^i\gamma_{\cC_0}(L_i).
\]
One shows that the element of $K(\cC_0)$ thus constructed is independent of the resolution $L$ chosen. If one has a second resolution $L'$, one ``caps'' $L$ and $L'$ by a third resolution $L''$, with surjective homomorphisms from $L''$ to $L$ and $L'$; this is possible by (i) and (ii). One is then reduced to proving that $E_{\cC_0}(L)=E_{\cC_0}(L'')$. But the difference of the two is $E_{\cC_0}(N)$, where $N$ is the kernel of the surjective homomorphism from $L''$ to $L$; this kernel is in $\cC_0$ by (i). Since $N$ is an acyclic complex in all dimensions, $E_{\cC_0}(N)=0$, whence the conclusion. Here $E_{\cC_0}(L)$, for a finite complex with values in $\cC_0$, denotes the Euler--Poincar\'e characteristic
\[
 \sum_i (-1)^i\gamma_{\cC_0}(L_i).
\]
One proves analogously that the function $\ell:\cC_1\to K(\cC_0)$ thus constructed is additive; hence it defines a homomorphism from $K(\cC_1)$ to $K(\cC_0)$, and it is immediate that this homomorphism is inverse to the canonical homomorphism from $K(\cC_0)$ to $K(\cC_1)$. QED.

\textbf{Remark.} To verify (ii), it is enough to verify that $\cC_0$ is contained in a subclass $\cC_2$ of $\cC$ satisfying the two conditions
\begin{enumerate}[label=(ii \alph*)]
\item every $A\in\cC_1$ is isomorphic to a quotient of an $L\in\cC_2$;
\item if $u:A\to B$, with $A,B\in\cC_2$, then $\Ker u\in\cC_1$.
\end{enumerate}

\textbf{Corollary.} Let $X$ be a quasi-projective algebraic space. Let $\cF(X)$ be the category of coherent algebraic sheaves on $X$, let $\xi(X)$ be the category of coherent locally free algebraic sheaves, equivalently vector bundles on $X$, and let $\cF_0(X)$ be the category of coherent algebraic sheaves $F$ such that, for every $x\in X$, the stalk $F_x$ is an $\cO_x$-module of finite cohomological dimension. Thus $\xi(X)\subset\cF_0(X)\subset\cF(X)$. Then the canonical homomorphism
\begin{equation}\tag{2.12}
 K(\xi(X))\longrightarrow K(\cF_0(X))
\end{equation}
is an isomorphism.

We verify the conditions of Theorem 2.3. Conditions (i) and (iii) follow from the definition of cohomological dimension and from the fact that a finite-type projective module over a noetherian local ring $\cO_x$ is free. To verify (ii), we shall prove that $\cC=\cF(X)$ satisfies conditions (ii a) and (ii b) of the remark. The second is trivial; therefore it remains to prove the following result.

\textbf{Proposition 2.4 (Serre).} Let $X$ be a quasi-projective algebraic space. Then every coherent algebraic sheaf on $X$ is isomorphic to the quotient of a locally free algebraic sheaf.

This proposition is true if $X$ is a closed subset of projective space by [2], since, for $n$ large enough, $F(n)$ is generated by its sections, which means that $F$ is a quotient of $\cO(-n)^k$. In the general case, $X$ is an open subset of a closed subset of a projective space, and it suffices to use the following result, due to Cartier and Serre in particular cases.

\textbf{Proposition 2.5.} Let $F$ be a coherent algebraic sheaf on an open subset $U$ of an algebraic space $X$. Then $F$ is isomorphic to the restriction of a coherent algebraic sheaf on $X$.

By applying Zorn's lemma to the opens containing $U$ on which one can extend $F$, one is reduced to the case where $X$ is affine. Let $\overline F$ be the algebraic sheaf on $X$ whose sections on an open $V$ are the sections of $F$ on $U\cap V$. One checks easily that $\overline F$ is quasi-coherent, i.e. defined on every affine open $V$---hence here on all of $X$---by a module, not necessarily of finite type, over the coordinate ring of that open. This is immediate if $U$ itself is affine, since $\overline F$ is then defined on the open $V$ by $\Gamma(U\cap V,F)$ considered as a module over $\Gamma(V,\cO_X)$. In the general case, $U$ is a finite union of affine opens $U_i$, hence $\overline F$ is the kernel of the evident homomorphism of quasi-coherent sheaves
\[
 \prod_i \overline F\big|_{U_i}\longrightarrow \prod_{i,j}\overline F\big|_{U_i\cap U_j},
\]
and is therefore itself quasi-coherent. Since $X$ is affine, $\overline F$ is generated by its sections. Since $\overline F|_U=F$ is coherent, there exists a finite number of sections of $\overline F$ which generate $F$ on $U$. These sections generate a coherent subsheaf of $\overline F$ whose restriction to $U$ is isomorphic to $F$. QED.

Now suppose that $X$ is nonsingular and quasi-projective. By the theorem on syzygies, one has, with the notation of the corollary to Theorem 2.3,
\[
 \cF_0(X)=\cF(X).
\]
In this case, therefore,
\begin{equation}\tag{2.13}
 K(\xi(X))=K(\cF(X)).
\end{equation}
One may express this isomorphism as follows.

\textbf{Corollary 2 to Theorem 2.3.} Let $X$ be a nonsingular quasi-projective variety. For every abelian group $A$, there is a one-to-one correspondence between additive functions $f$ of vector bundles on $X$ with values in $A$, and additive functions $g$ of coherent algebraic sheaves on $X$ with values in $A$. To every $f$ there corresponds the unique $g$ such that
\begin{equation}\tag{2.14}
 f(E)=g(\cO_X(E))
\end{equation}
for every vector bundle $E$ on $X$, where $\cO_X(E)$ denotes the sheaf of germs of regular sections of $E$ on $X$.

To compute $g(F)$ for an arbitrary coherent algebraic sheaf, one considers a finite resolution of $F$ by sheaves of the form $\cO_X(E_i)$, and one then has
\begin{equation}\tag{2.14 bis}
 g(F)=\sum_i (-1)^i f(E_i).
\end{equation}

In particular, if $X$ is connected, in \S 1 we considered an additive function $\Ctil(E)$ of the vector bundle $E$, with values in the group $A=\Atil(X)$. We deduce from it an additive function of coherent sheaves, say $\Ctil(F)$, and, by passing to components, functions
\[
 F\longmapsto c^i(F)\in A^i(X).
\]
The $c^i(F)$ are called the \emph{Chern classes of the coherent sheaf} $F$. Theoretically they are computed from the Chern classes of vector bundles by (2.14 bis), with the understanding that in the second member one must pass to multiplicative notation. As for the component of $\Ctil(F)$ in $\Z$, it is the rank of the sheaf $F$, corresponding to the rank of a module over an integral domain: if $X$ is affine, to which one can always reduce, $F$ is defined by a finite-type module over the coordinate ring of $X$, and one takes the rank of this module.

\subsection*{3. Functorial generalities on $K(X)$}

If $X$ is an algebraic space, put, for brevity,
\[
 K(X)=K(\cF(X)).
\]
When $X$ is nonsingular and quasi-projective, identify this group with $K(\xi(X))$ by the corollary 2 to Theorem 2.3. Since $K(\xi(X))$ is not only an abelian group but a $\lambda$-ring, it follows by transport of structure that $K(X)$ is also a $\lambda$-ring, and is therefore provided with a multiplicative structure and operations $\lambda^i$. It is important to define these structures directly, without passing through vector bundles; this will be done below.

If $E$ is a vector bundle on $X$, denote by $\gamma(E)$ its class in $K(X)$. If $F$ is a coherent algebraic sheaf on $X$, $\gamma(F)$ denotes in the same way its class in $K(X)$. If $Y$ is an irreducible subvariety of $X$, put $\gamma(Y)=\gamma(\cO_Y)$, where $\cO_Y$ denotes the sheaf of local rings of $Y$, considered as a coherent algebraic sheaf on $X$. By linearity, one defines the symbol $\gamma(Y)$ when $Y$ is any cycle on $X$. When there is risk of confusion, one writes $\gamma_X$ for $\gamma$.

Let $X$ be an arbitrary algebraic space. We introduce on $K(X)$ a filtration, denoting by $K(X)^i$ the subgroup of $K(X)$ generated by the $\gamma(F)$ with
\[
 \dim \Supp F\leq n-i,
\]
where $n$ is the dimension of $X$. From Serre's d\'evissage lemma, in the form given in [3], follows the following result.

\textbf{Proposition 2.6.} The group $K(X)^i$ is generated by the $\gamma(Y)$, where $Y$ runs through the subvarieties of dimension $n-i$ of $X$. If $X$ is irreducible, then $\Gr^0(K(X))=\Z$.

The latter isomorphism is given by the homomorphism from $K(X)$ to $\Z$ defined with the aid of the notion of rank of a coherent algebraic sheaf.

Now suppose that $X$ is nonsingular. Following the developments at the beginning of \S 2, define a multiplication in $K(X)$ by putting
\begin{equation}\tag{2.15}
 \gamma(F)\gamma(G)=\sum_{i\geq 0}(-1)^i\gamma(\Tor_i(F,G)),
\end{equation}
the $\Tor_i$ being taken, of course, with respect to the local rings of $X$. The exact sequence of the $\Tor_i$ shows that formula (2.15) indeed defines a biadditive map from $K(X)\times K(X)$ to $K(X)$; associativity follows from the spectral sequence of multiple Tors. One also verifies in this way that one can compute the product of several factors by the formula
\begin{equation}\tag{2.16}
 \gamma(F_1)\cdots\gamma(F_n)=\sum_{i\geq 0}(-1)^i\gamma(\Tor_i(F_1,\ldots,F_n)),
\end{equation}
where, in the second member, the $\Tor_i$ are simultaneous Tors, computed by simultaneous projective resolution of the arguments $F_i$.

If in (2.15) one supposes that $G$ is locally free, one finds
\begin{equation}\tag{2.15 bis}
 \gamma(F)\gamma(G)=\gamma(F\otimes G),
\end{equation}
which shows in particular that $\gamma(\cO_X)$ is the unit element of $K(X)$, and that the canonical homomorphism from $K(\xi(X))$ to $K(X)$ is a homomorphism of rings.

Formula (2.15) is evidently suggested by Serre's ``Tor formula'' in intersection theory. This formula proves the second part of the following proposition, the first part following immediately from (2.15) and from a local Tor calculation.

\textbf{Proposition 2.7.} Let $X$ be a nonsingular variety, and let $Y$ and $Z$ be two cycles in $X$. If $Y$ and $Z$ meet everywhere transversely, then
\begin{equation}\tag{2.17}
 \gamma(Y)\gamma(Z)=\gamma(Y\cdot Z).
\end{equation}
If $Y$ and $Z$ are homogeneous and of respective dimensions $n-i$ and $n-j$, and if they meet without an excessive component, then
\begin{equation}\tag{2.18}
 \gamma_X(Y)\gamma_X(Z)=\gamma_X(Y\cdot Z)\mod K(X)^{i+j+1}.
\end{equation}
One must be careful: equality does not in general hold in (2.18). We shall see in \S 4 that, if $X$ is nonsingular and quasi-projective, the ring structure of $K(X)$ is compatible with its filtration, and that the associated graded ring $\Gr(K(X))$ is identified with a quotient of $A(X)$ by a subgroup of torsion, perhaps always zero.\footnote{No: for a counterexample see this seminar, XIV 4.7.}

I do not know how to define directly the $\lambda^i(F)$ in terms of sheaves except when the field $k$ has characteristic $0$; this is the reason that has prevented us from proving the Riemann--Roch theorem when $k$ is not of characteristic $0$. Let $F$ be a coherent algebraic sheaf on $X$ on which the group $\mathfrak S_p$ of permutations of $p$ letters operates. Denote by $F^{\mathrm{alt}}$ the set of $x\in F$ such that
\[
 s\cdot x=\varepsilon(s)x
\]
for every $s\in\mathfrak S_p$, where $\varepsilon(s)$ is the signature of $s$. It is evidently a coherent algebraic subsheaf of $F$. This said, one has:

\textbf{Theorem 2.8.} Let $X$ be a nonsingular quasi-projective variety over a field $k$ of characteristic $0$. Let $F$ be a coherent algebraic sheaf on $X$. Then
\begin{equation}\tag{2.19}
 \lambda^p(\gamma(F))=\sum_{i\geq 0}(-1)^i\gamma\left((\Tor_i(F,F,\ldots,F))^{\mathrm{alt}}\right).
\end{equation}
Here it is clear that the group $\mathfrak S_p$ operates on the sheaf $\Tor_i(F,\ldots,F)$.

\emph{Principle of the proof.} One considers a finite resolution $L$ of $F$ by locally free sheaves $L_i$, and uses the following general lemma.

\textbf{Lemma 2.9.} Let $L=(L_i)$ be a graded locally free coherent algebraic sheaf with only finitely many nonzero components. In $K(X)$ one has the formula
\begin{equation}\tag{2.20}
 E\left((\otimes^p L)^{\mathrm{alt}}\right)=\lambda^p(E(L)).
\end{equation}
In this formula, $\otimes^pL$ is considered as a sheaf on which the group $\mathfrak S_p$ acts, taking the gradings into account; the gradings therefore introduce signs in the usual way, cf. Cartan--Eilenberg. Moreover, if $M$ is a graded coherent algebraic sheaf on $X$, put
\[
 E(M)=\sum_i(-1)^i\gamma(M_i).
\]

To prove formula (2.20), one is reduced to proving the analogous formula in $K(G)$, where $G$ is an algebraic group, for instance a product of groups $\operatorname{Gl}(n,k)$, and where $L$ is a finite-dimensional $G$-module over $k$ (cf. Chapter I, \S 2, Example 2, for the definition of $K(G)$). In that case, this formula has already been used in Theorem 1.4 to show that the second member of (1.33) satisfies the same functional equation as the first member; the verification is elementary. QED.

Let $f$ be a morphism from $X$ to $Y$, where $Y$ is an algebraic space. The notion of inverse image of vector bundle defines a homomorphism
\[
 f^!:K(\xi(Y))\longrightarrow K(\xi(X))
\]
of $\lambda$-rings. Identifying vector bundles with locally free coherent algebraic sheaves, the notion of inverse image of vector bundle corresponds to the notion of inverse image of algebraic sheaf, already used in [3, no. 8]. When $Y$ is nonsingular, one defines more generally a homomorphism of groups
\begin{equation}\tag{2.21}
 f^!:K(Y)\longrightarrow K(X)
\end{equation}
by a formula involving an alternating sum of Tors, inspired by the developments at the beginning of \S 2, and which the reader will spell out. One then has commutativity in the following diagram:
\[
\begin{tikzcd}[column sep=large,row sep=large]
K(\xi(Y)) \arrow[r,"f^!"] \arrow[d] & K(\xi(X)) \arrow[d]\\
K(Y) \arrow[r,"f^!"'] & K(X),
\end{tikzcd}
\]
which shows in particular, when $X$ and $Y$ are both nonsingular and quasi-projective, that the two definitions of $f^!$ coincide under the identifications of $K(X)$ with $K(\xi(X))$ and of $K(Y)$ with $K(\xi(Y))$.

Another case where one can define a homomorphism (2.21), without $Y$ necessarily nonsingular, is the case where the functor $F\mapsto f^!(F)$ from $\cF(Y)$ to $\cF(X)$ is exact, i.e. where the local rings of $X$ are flat modules over the local rings of $Y$. It is enough then to put
\[
 f^!(\gamma_Y(F))=\gamma_X(f^!(F)).
\]
This case occurs in particular when $X$ is a locally trivial fibre space over $Y$; to see this, one is reduced to the case of a product $X=Y\times T$. Note that one also verifies that in this case one has $f^!(\cO_Z)=\cO_{f^{-1}(Z)}$ if $Z$ is an irreducible subvariety of $Y$, whence
\begin{equation}\tag{2.22}
 f^!(\gamma_Y(Z))=\gamma_X(f^{-1}(Z))
\end{equation}
if $X$ is flat over $Y$. Without this latter condition, formula (2.22) becomes false in general, but one has the following analogue of Proposition 2.7.

\textbf{Proposition 2.8.} Let $f:X\to Y$ be a morphism of nonsingular varieties, and let $Z$ be a cycle on $Y$. If $Z$ is everywhere transverse to $f$, formula (2.22) is exact. If $Z$ is homogeneous of codimension $i$ and if $f^{-1}(Z)$ has no excessive component, then
\begin{equation}\tag{2.23}
 f^!(\gamma_Y(Z))=\gamma_X(f^{-1}(Z))\mod K(X)^{i+1}.
\end{equation}
The proof is analogous to that of Proposition 2.7.

Finally, let $X$ and $Y$ again be arbitrary algebraic spaces, and let $f$ be a proper morphism from $X$ to $Y$. For every coherent algebraic sheaf $F$ on $X$, the direct image sheaf $f_*(F)$ is coherent algebraic, and the same holds more generally for the sheaves $R^if_*(F)$ (cf. [3]). Recall that, by definition, $R^if_*(F)$ is the sheaf on $Y$ associated with the presheaf which sends an open $U$ to the group $H^i(f^{-1}(U),F)$; in particular
\begin{equation}\tag{2.24}
 R^0f_*(F)=f_*(F),\qquad R^qf_*(F)=0\quad(q>\dim X),
\end{equation}
and moreover
\begin{equation}\tag{2.25}
 \Gamma(U,R^qf_*(F))=H^q(f^{-1}(U),F)
\end{equation}
if the open $U$ is affine. According to \S 2, one defines a homomorphism from $K(X)$ to $K(Y)$, still denoted $f_*$, by the formula
\begin{equation}\tag{2.26}
 f_*(\gamma(F))=\sum_{q=0}^{\infty}(-1)^q\gamma(R^qf_*(F)).
\end{equation}
Thus $K(X)$ becomes a covariant functor in $X$ relative to proper morphisms. The relation $(fg)_*=f_*g_*$, just like the relation $(fg)^!=g^!f^!$ which we forgot to mention in its place, is a particular case of the transitivity relations alluded to in \S 2 and which follow from the spectral sequences of composed functors.

Let us note the following formula, valid for $Y$ nonsingular and particularly easy to verify when $Y$ is also quasi-projective by taking for $y$ the class of a vector bundle on $Y$:
\begin{equation}\tag{2.27}
 f_*(x\cdot f^!(y))=f_*(x)\cdot y\qquad (x\in K(X),\ y\in K(Y)).
\end{equation}
In the particular case where the inverse images of the points of $Y$ by $f$ are finite and contained in an affine open,\footnote{This latter condition is evidently a consequence of the first, something the author still did not know in 1957!} one checks that the functor $F\mapsto f_*(F)$ is exact and that $R^qf_*(F)=0$ for $q>0$. In this case, therefore, $f_*(\gamma_X(F))=\gamma_Y(f_*(F))$. In particular, if $X$ is an algebraic subspace of $Y$ and if $i$ is the canonical injection of $X$ into $Y$, then
\begin{equation}\tag{2.28}
 i_*(\gamma_X(F))=\gamma_Y(F),
\end{equation}
where, in the second member, $F$ is considered as a sheaf on $Y$ which is zero outside $X$.

\subsection*{4. Some technical results}

Let $X$ be an algebraic space. We have seen in \S 3 that the projection $p$ from $X\times k$ to $X$ defines a homomorphism
\begin{equation}\tag{2.29}
 p^!:K(X)\longrightarrow K(X\times k).
\end{equation}

\textbf{Proposition 2.9.} The homomorphism (2.29) is surjective.

We argue by induction on $n=\dim X$. The assertion is trivial if $n<0$; suppose $n\geq 0$. By Proposition 2.6, it is enough to verify that for every irreducible closed subset $Y$ of $X\times k$, $\gamma(Y)$ is in the image of $p^!$. It is enough to continue the reasoning by supposing $X$ irreducible. If $p(Y)$ is not dense in $X$, our assertion follows from the induction hypothesis. Otherwise, either $Y=X\times k$ and the assertion is trivial, or $Y$ is a hypersurface in $X\times k$. Let $U$ be an affine open of $X$ formed of nonsingular points. If $D$ is the divisor on $U\times k$ induced by $Y$, one knows by commutative algebra that it is linearly equivalent to a divisor of the form $D'\times k$, where $D'$ is a divisor on $U$. Then
\[
 \gamma_{U\times k}(D)-\gamma_{U\times k}(D'\times k)
\]
belongs to $K(U\times k)^2$ by the corollary to Proposition 2.10, and by what has already been proved this element is of the form $p^!(x')$ with $x'\in K(U)$. Thus
\[
 \gamma_{U\times k}(D)=p^!(x'+\gamma_U(D')),
\]
but $x'$ is the restriction of an element $y$ of $K(X)$ by Proposition 2.10, and consequently the restriction of $\gamma_{X\times k}(D)-p^!(y)$ to $U\times k$ is zero. By Proposition 2.10, $\gamma_{X\times k}(D)-p^!(y)$ comes from an element of
\[
 K(X\times k-U\times k)=K(X'\times k),\qquad X'=X-U.
\]
The proof is finished by applying the induction hypothesis again. QED.

\textbf{Corollary.} Let $X$ be an algebraic space and let $p$ be the projection from $X\times k^n$ to $X$. Then the homomorphism
\[
 p^!:K(X)\longrightarrow K(X\times k^n)
\]
is surjective; it is even bijective if $X$ is nonsingular.

The first assertion is by induction on $n$. For the second, put $f(x)=(x,0)$, whence $pf=1$ and $f^!p^!=1$, which proves that $p^!$ is injective.

\textbf{Proposition 2.10.} Let $X$ be an algebraic space, let $U$ be an open subset of $X$, and let $Y$ be its complement in $X$. Then the sequence of natural homomorphisms
\begin{equation}\tag{2.30}
 K(Y)\longrightarrow K(X)\longrightarrow K(U)\longrightarrow 0
\end{equation}
is exact.

N.B. Compare with Proposition 2.3.

The fact that $K(X)\to K(U)$ is surjective follows at once from Proposition 2.5, or also from Proposition 2.6. It remains to prove that if $x\in K(X)$ has zero restriction in $K(U)$, then it is in the image of $K(Y)$. For this, return to the explicit definition of $K(X)$ and $K(U)$. There are two coherent algebraic sheaves $F$ and $G$ on $X$ with
\[
 x=\gamma_X(F)-\gamma_X(G)
\]
and
\[
 \gamma_U(F)-\gamma_U(G)=0.
\]
This latter relation means that one can find coherent sheaves $P'$ and $Q'$ on $U$, provided with filtrations whose factors are pairwise isomorphic on $U$, and such that one has an isomorphism on $U$
\[
 P=F+Q'\longrightarrow Q=G+P'.
\]
By Proposition 2.5, the sheaves $P'$ and $Q'$ are restrictions to $U$ of coherent sheaves on $X$, denoted by the same letters; in the same way, their filtrations come from filtrations on the corresponding sheaves on $X$, by Lemma 2.11 below. One then has
\[
 \gamma_X(F)-\gamma_X(G)=(\gamma_X(P)-\gamma_X(Q))+\bigl(\gamma_X(P')-\gamma_X(Q')\bigr).
\]
Introducing the factors $P_i$ and $Q_i$ of the filtrations of $P'$ and $Q'$ on $X$, the last term is the sum of terms $\gamma_X(P_i)-\gamma_X(Q_i)$, where $P_i|_U\simeq Q_i|_U$. We are thus reduced to proving the following particular case: if $R$ and $S$ are coherent sheaves on $X$ such that $R|_U\simeq S|_U$, then $\gamma_X(R)-\gamma_X(S)$ is in the image of $K(Y)$. Let $T$ be the graph of the given isomorphism from $R|_U$ to $S|_U$; it is a coherent subsheaf of $(R\times S)|_U$, hence by Lemma 2.11 it is the restriction to $U$ of a coherent subsheaf, again denoted $T$, of $R\times S$. Then
\[
 \gamma_X(R)-\gamma_X(S)=(\gamma_X(R)-\gamma_X(T))-(\gamma_X(S)-\gamma_X(T)),
\]
and it is enough to prove, for instance, that the first term of the second member is in the image of $K(Y)$. But the projection $u:T\to R$ is an isomorphism over $U$, hence
\[
 \gamma_X(R)-\gamma_X(T)=\gamma_X(\Coker u)-\gamma_X(\Ker u),
\]
and $\Ker u$ and $\Coker u$ have support in $Y$.

It remains, then, to prove the following lemma.

\textbf{Lemma 2.11.} Let $X$ be an algebraic space, let $F$ be a coherent algebraic sheaf on $X$, let $U$ be an open subset of $X$, and let $G$ be a coherent subsheaf of $F|_U$. Then $G$ is the restriction to $U$ of a coherent subsheaf of $F$.

Let $\overline G$ be the subsheaf of $F$ whose sections on an open $V$ of $X$ are the sections of $F$ on $V$ whose restriction to $U\cap V$ is a section of $G$. Everything amounts to proving that this sheaf $\overline G$ is coherent. For this one may suppose $X$ and $U$ affine, because $U$ is a union of affine opens $U_i$ and $\overline G$ is then the intersection of the analogous sheaves $\overline{G|U_i}$. Then the coordinate ring of $U$ is a ring of fractions $A_S$ of the coordinate ring $A$ of $X$, and one immediately verifies that if $F$ is defined by the $A$-module $M$, then $F|_U$ is defined by the $A_S$-module
\[
 M_S=A_S\otimes_A M.
\]
The subsheaf $G$ of $F|_U$ is then defined by a submodule $N'$ of $M_S$, and one sees that $\overline G$ is precisely defined by the submodule $N$ of $M$, inverse image of $N'$ under the canonical homomorphism from $M$ to $M_S$. QED.

\textbf{Corollary 1 to Proposition 2.10.} Let $x\in K(X)$. Then $x\in K(X)^i$ if and only if there exists a closed subset $Y$ of $X$, of dimension $\leq n-i$, such that the restriction of $x$ to $X-Y$ is zero.

\textbf{Corollary 2.} Suppose $X$ is connected of dimension $n$, with no singularities in dimension $n-1$. If $D$ and $D'$ are two linearly equivalent divisors on $X$, then
\[
 \gamma(D)=\gamma(D')\mod K(X)^2.
\]
By Corollary 1, one may suppose $X$ nonsingular. The corollary then follows from the general formula
\begin{equation}\tag{2.31}
 \gamma(D)=1-\gamma(L(-D))\mod K(X)^2,
\end{equation}
where $L(-D)$ denotes the rank-one vector bundle defined by the divisor $-D$. To verify (2.31), one writes it in the form
\begin{equation}\tag{2.31 bis}
 \gamma(L(-D))=1-\gamma(D)\mod K(X)^2,
\end{equation}
and notes that the two members are homomorphisms from the group of divisors to the multiplicative group of invertible elements of the ring $K(X)/K(X)^2$. Here $K(X)^1/K(X)^2$ has square zero, which is very easily verified on the explicit formula (2.15), using Corollary 1. Thus one is reduced to the case where $D$ is an irreducible hypersurface. More generally, if $D$ is a hypersurface whose components are disjoint, one has the equality
\begin{equation}\tag{2.32}
 \gamma(D)=1-\gamma(L(-D)),
\end{equation}
as follows from the exact sequence of sheaves
\[
0\longrightarrow \cO_X(L(-D))\longrightarrow \cO_X\longrightarrow \cO_D\longrightarrow 0.
\]

\textbf{Theorem 2.12.} Let $X$ be a nonsingular quasi-projective variety, and let $Z$ and $Z'$ be two homogeneous cycles of codimension $p$ on $X$ which are linearly equivalent. Then
\begin{equation}\tag{2.33}
 \gamma_X(Z)=\gamma_X(Z')\mod K(X)^{p+1}.
\end{equation}
By hypothesis, one can find a cycle $Y$ of dimension $n-p+1$ on $X\times k$ ($n$ being the dimension of $X$), and two points $a,a'\in k$, such that
\[
 Z=Y\cdot (X\times a),\qquad Z'=Y\cdot (X\times a').
\]
By formula (2.23), one then has
\[
 \gamma_X(Z)=i_a^!(\gamma_{X\times k}(Y))\mod K(X)^{p+1},
\]
and an analogous formula for $\gamma_X(Z')$. It is therefore enough to prove that $i_a^!=i_{a'}^!$, which follows at once from Proposition 2.9. QED.

\textbf{Corollary 1.} The filtration of $K(X)$ is compatible with its ring structure, i.e. one has
\[
 K(X)^iK(X)^j\subset K(X)^{i+j}
\]
for all integers $i,j\geq 0$. Moreover there is a natural surjective homomorphism of graded rings
\begin{equation}\tag{2.34}
 \varphi:A(X)\longrightarrow \Gr(K(X)),
\end{equation}
where the second member denotes the graded ring associated with $K(X)$, provided with the filtration $(K(X)^i)_{i\geq 0}$, defined by the condition
\[
 \varphi(\ellc(Z^{n-p}))=\gamma(Z^{n-p})\mod K(X)^{p+1}.
\]
This corollary follows immediately from Theorem 2.12, Propositions 2.6 and 2.7, and Chow's theorem 2.1, by induction on the integer $i+j$.

\textbf{Corollary 2.} Let $X$ and $Y$ be two nonsingular quasi-projective varieties and let $f$ be a morphism from $X$ to $Y$. Then the homomorphism $f^!$ from $K(Y)$ to $K(X)$ is compatible with the filtrations, and one has a commutative diagram
\[
\begin{tikzcd}[column sep=large,row sep=large]
A(Y) \arrow[r,"\varphi_Y"] \arrow[d,"f^*"'] & \Gr(K(Y)) \arrow[d,"\Gr(f^!)"]\\
A(X) \arrow[r,"\varphi_X"'] & \Gr(K(X)).
\end{tikzcd}
\]
Using the graph of $f$, decompose $f$ as the product of an immersion of $X$ into $X\times Y$ and of a projection of $X\times Y$ onto $Y$. For the latter the corollary is trivial, and for an immersion one proceeds as for Corollary 1, using Proposition 2.8 in place of Proposition 2.7.

\textbf{Proposition 2.13 (``cellular decomposition'').} Let $P$ be an algebraic space, irreducible or not, provided with an increasing family
\[
P^0\subset P^1\subset\cdots\subset P^n=P
\]
of closed subsets, such that for every $i$ with $0\leq i\leq n$, the connected components of $P^i-P^{i-1}$ are algebraic spaces $V_{i,\alpha}$ isomorphic to $k^i$; put $P^{-1}=\varnothing$. Then:
\begin{enumerate}[label=\alph*)]
\item For every algebraic space $X$, the natural map from $K(X)\otimes K(P)$ to $K(X\times P)$, deduced from the tensor product of sheaves on $X$ and $P$, is surjective.
\item The $\gamma_P(V_{i,\alpha})$ form a system of generators of the group $K(P)$.
\end{enumerate}
The proof is by induction on $n$, using the corollary to Proposition 2.9 and Proposition 2.10. N.B. Unfortunately, I do not know whether the $\gamma_P(V_{i,\alpha})$ form a basis of $K(P)$ or whether the homomorphism considered in a) is bijective. We shall not need to know this below.

\textbf{Theorem 2.14.} Let $x_i$ ($1\leq i\leq q$) and $E_j$ ($1\leq j\leq r$) be letters, and let $n_j$ ($1\leq j\leq r$) be positive integers. Consider the symbols $\lambda^k(x_i)$ with $k\geq 1$ and $\lambda^k(E_j)$ with $1\leq k\leq n_j$ as indeterminates, and let $R$ be a polynomial with integral coefficients in these indeterminates. Suppose that, for every graded ring $A$, when one substitutes in $R$ elements $\xi_i$ for the $x_i$ and elements
\[
 \eta_j=[n_j,1+\eta_j^{(1)}+\cdots+\eta_j^{(n_j)}]
\]
of $\Atil$ for the $E_j$, subject to the restrictions $\varepsilon(\xi_i)=0$ and $\eta_j^{(k)}=0$ for $k>n_j$, one has
\[
 R\bigl((\lambda^k(\xi_i)),(\lambda^k(\eta_j))\bigr)\in \Atil^p.
\]
Then, if $X$ is an irreducible nonsingular quasi-projective variety, for arbitrary elements $x'_i$ of $K(X)^1$ and vector bundles $E'_j$ of dimension $n_j$ on $X$ ($1\leq i\leq q$, $1\leq j\leq r$), one has
\begin{equation}\tag{*}
 R\bigl((\lambda^k(x'_i)),(\lambda^k(\gamma(E'_j)))\bigr)\in K(X)^p.
\end{equation}

\emph{Proof.} By \S 2, the $x'_i$ are themselves differences of elements of the form $\gamma(E)$, with $E$ a vector bundle. On the other hand, assuming $X$ locally closed in a projective space $P$, it follows from \S 2 and Proposition 2.5 that, for every $F\in\cF(X)$, the sheaf $F(n)$ is generated by its sections for $n$ sufficiently large. If $F$ is defined by a vector bundle $E$, this proves that $E(n)$ is the inverse image of the canonical vector bundle on a Grassmannian variety $G$ by a morphism from $X$ to $G$. Considering the corresponding morphism from $X$ to $P\times G$, one sees that $F$ is the inverse image of a vector bundle on $P\times G$. Applying this to all vector bundles which occur, and using Corollary 2 to Theorem 2.12, we are reduced to the case where $X$ is a product of Grassmannians. In this case the conclusion will follow from the following lemma.

\textbf{Lemma 2.15.} Suppose that $X$ is a product of Grassmannians. Then one can find a $\lambda$-homomorphism from $K(X)$ to a torsion-free $\lambda$-ring of the form $\Atil$, where $A$ is a graded ring, compatible with filtrations, and such that the corresponding homomorphism $\Gr(K(X))\to\Gr(A)$ is injective.

To prove the lemma, we shall use Chow's theory of Chern classes. If one could prove the lemma, or Theorem 2.14 directly, one would deduce a purely ``sheaf-theoretic'' theory of Chern classes; cf. \S 5.

We shall consider the homomorphism $\Ctil$ from $K(X)$ to $\Atil(X)$ and put $A=A(X)$. It remains to prove that the homomorphism $\psi=\Gr(\Ctil)$ from $\Gr(K(X))$ to $\Gr(A(X))\simeq A(X)$ is injective. The fact that $\Ctil$ is compatible with filtrations is true for every nonsingular quasi-projective variety and follows easily from the corollary to Proposition 2.3. On the other hand, since $X$ is a product of Grassmannians $G_j$, there are vector bundles $E_j$ on $X$, inverse images of the canonical bundles on the factors. We shall use the following known facts:
\begin{enumerate}[label=\alph*)]
\item On $X$, linear equivalence of cycles is equivalent to numerical equivalence, and $A(X)$ is a group of finite rank.
\item The $c^i(E_j)$ generate the ring $A(X)$.
\end{enumerate}
From a) it follows that
\[
 \varphi:A(X)\longrightarrow \Gr(K(X))
\]
is injective, and hence bijective, by the following lemma.

\textbf{Lemma 2.15.} Let $X^n$ be a nonsingular projective variety and let $Z^{n-p}$ be a cycle on $X$ such that $\varphi(Z)\in\Gr^p(K(X))$ is zero. Then $Z^{n-p}$ is numerically equivalent to zero.

To prove this lemma, it suffices to use the additive function
\[
 \chi(F)=\sum_{i\geq 0}(-1)^i\dim H^i(X,F)
\]
and the fact that this function is equal to $1$ for $F=\cO_{(x)}$ for every $x\in X$.

Thus, in the case under consideration, $A(X)$ and $\Gr(K(X))$ have the same finite rank in every degree, and $\Gr(K(X))$ is torsion-free, since this is true of $A(X)$. From assertion b) above it follows, in a formal way, that the map $\Ctil$ from $K(X)$ to $\Atil(X)$ is surjective modulo finite groups. Since $\Gr(K(X))$ and $A(X)$ have the same finite rank in every degree, it then follows easily that the homomorphism $\Gr(\Ctil)$ from $\Gr(K(X))$ to $\Gr(\Atil(X))$ is bijective modulo finite groups in every degree; one uses induction on the degree. Since $\Gr(K(X))$ is torsion-free, this homomorphism $\Gr(\Ctil)$ is therefore injective. QED.

N.B. The homomorphism $\Ctil$ is not surjective in general, even if $X$ is a projective space of dimension $2$.

\subsection*{5. Sheaf-theoretic definition of Chern classes. Application to the study of injection morphisms}

\textbf{Theorem 2.16.} Let $X$ be a nonsingular quasi-projective variety. Then, for every $x\in K(X)$ and every integer $i\geq 0$, the element $\gamma^i(x-\varepsilon(x))$ is of filtration $\geq i$. Let $c^i(x)\in\Gr^i(K(X))$ be the class of this element modulo $K(X)^{i+1}$. If one puts
\[
 \widetilde c(x)=\left[\varepsilon(x),1+\sum_{i\geq 0}c^i(x)\right]\in \widetilde{\Gr(K(X))},
\]
then $\widetilde c$ satisfies conditions analogous to those of Theorem 2.2.

The first assertion follows from Theorem 2.14, as does the fact that $x\mapsto\widetilde c(x)$ is a homomorphism of $\lambda$-rings; formula (1.25) implies $\gamma_t(x+y)=\gamma_t(x)\gamma_t(y)$, which already shows that $\widetilde c$ is an additive homomorphism. The functoriality of $\widetilde c$ is verified immediately, and (2.8) follows from (2.31). QED.

\textbf{Corollary 1.} One has
\[
 c^i(x)=\varphi(c^i(x)).
\]
This follows from a uniqueness theorem for Chern classes, proved as in Hirzebruch by passage to flag varieties. One only has to show that, if $Y$ is a flag variety associated with a vector bundle on $X$, then $\Gr(K(X))\to\Gr(K(Y))$ is injective; this is verified by the explicit determination of $K(Y)$, with all its structures, from $K(X)$. I shall not give the details here.

\textbf{Corollary 2.} Suppose that for every element $Z\in A(X)$ there exists a morphism $f$ from $X$ to a product $X'$ of Grassmannians such that $Z$ is in the image of $f^*$. Then the canonical homomorphisms
\[
 \varphi:A(X)\longrightarrow \Gr(K(X)),\qquad
 \psi:\Gr(K(X))\longrightarrow A(X)
\]
satisfy, in dimension $i$, the relations
\begin{equation}\tag{2.35}
 \psi\varphi=(-1)^{i-1}(i-1)!\,1,
 \qquad
 \varphi\psi=(-1)^{i-1}(i-1)!\,1.
\end{equation}
Since $\varphi$ is surjective, it is clearly enough to prove the first formula, and for this one may suppose that $X$ is a product of Grassmannians. Since then $\varphi$ is injective, it is enough to prove the second formula (2.35), which by Corollary 1 is written
\[
 \gamma^i(x-\varepsilon(x))=(-1)^{i-1}(i-1)!\,x\mod K(X)^{i+1}
\]
for every $x\in K(X)^i$. Since $\psi$ is injective, it is enough to prove the analogous relation for $\Ctil(x)\in A(X)^i$, and this is then a consequence of (1.22 bis).

\textbf{Remark.} Corollary 2 makes plausible the validity of (2.35) in all cases; there should exist a direct proof of it. In every case one proves easily the second formula (2.35) modulo the torsion subgroup of $\Gr(K(X))$, and in characteristic $0$ this same formula is a consequence of the Riemann--Roch theorem. The interest of (2.35), when true, lies in the fact that it shows that $\varphi$ is injective modulo torsion groups, perhaps always zero. Thus one does not lose much by defining Chern classes as elements of $\Gr(K(X))$. If one tensors with $\Q$, as will be done later for the Riemann--Roch theorem, the two definitions of Chern classes would then be identical.

Let $X$ still be nonsingular and quasi-projective. Let in addition $Y$ be a nonsingular subvariety of $X$, and let $i$ be the injection of $Y$ into $X$. For $i_*$ and $i^!$ one then has the following formulas:
\begin{enumerate}[label=(\roman*)]
\item $i^!$ is a homomorphism of $\lambda$-rings.
\item $i_*(x\cdot i^!(y))=i_*(x)\cdot y$, whence in particular
\[
 i_*i^!(y)=\xi y,
\qquad \text{with }\xi=i_*i^!(1)=\gamma_X(N).
\]
\item $i^!i_*(x)=x\cdot\lambda_{-1}(N)$, whence in particular
\[
 i^!(\xi)=i^!i_*(1)=\lambda_{-1}(N).
\]
\item $i_*(x)i_*(x')=i_*(xx'\lambda_{-1}(N))$.
\end{enumerate}
Here, by definition, $N$ is the class in $K(Y)$ of the conormal bundle of $Y$ in $X$. Formulas (i) and (ii) are true for arbitrary morphisms, not necessarily injections, and (iii) and (iv) are verified from the definition (2.15) and the analogous formula for inverse images, taking $x$ and $x'$ to be classes of vector bundles and using the fact that
\begin{equation}\tag{2.36}
 \Tor(\cO_Y,\cO_Y)=\Lambda(N).
\end{equation}
Formula (iv) expresses that $i_*$ is a homomorphism of rings if one introduces on $K(Y)$ a new multiplicative structure, by the definition
\[
 x\cdot_N x'=xx'\lambda_{-1}(N)
\]
(cf. Chapter I, \S 4). It is completed by the formula
\begin{enumerate}[label=(\roman*)]
\setcounter{enumi}{4}
\item $i_*(\lambda^p(N,x))=\lambda^p(i_*(x))$ for $p\geq 1$,
\end{enumerate}
which means that $i_*$ defines a $\lambda$-homomorphism from $K(Y)_N$ into $K(X)$. This formula (v), certainly true in all cases, is for the moment proved only in characteristic $0$, because it follows from the conjunction of Theorems 1.4 and 2.8.

There are analogous formulas for $i_*$ and $i^*$:
\begin{enumerate}[label=(\roman* bis)]
\item $i^*$ is a homomorphism of graded rings.
\item $i_*(x\cdot i^*(y))=i_*(x)y$, whence in particular
\[
 i_*i^*(y)=\xi' y,
\qquad \text{with }\xi'=i_*(1).
\]
\item $i^*i_*(x)=x\cdot (-1)^qN^q$, whence in particular
\[
 i^*(\xi')=(-1)^qN^q.
\]
\item $i_*(x)i_*(x')=i_*(xx'(-1)^qN^q)$.
\item $i_*$ increases degrees by $q$ units.
\end{enumerate}
In these formulas, $q$ is the codimension of $Y$ in $X$, and one puts $N^i=c^i(N)$. Only formulas (iii bis) and (iv bis) require proof. I do not know whether they are well known, or even whether they are true, and I confine myself to noting that they follow by reduction from formulas (ii) and (iv), provided one works in $\Gr(K(Y))$ and $\Gr(K(X))$ and not in $A(Y)$ and $A(X)$. One then uses the fact that $i_*$ and $i^*$ come by passage to the quotient from $i_*$ and $i^!$, noting that $i_*$ shifts the filtration by $q$ units.

All universal formulas relating $i_*$ and $i^!$, respectively $i_*$ and $i^*$, are formal consequences of the preceding formulas, as one can convince oneself without great difficulty.

Formula (v) gives
\[
 i_*(\gamma^p(N,x))=\gamma^p(i_*(x)).
\]
The conjunction of Theorem 2.14 and Proposition 1.5 shows that $\gamma^p(N,x)$, for $p=q+j$, is of filtration $\geq j$. If one denotes by $G_{q,j}(N,x)$ its class in $\Gr^j(K(Y))$, Theorem 2.16 shows that
\[
 i_*(G_{q,j}(N,x))=c^p(i_*(x)).
\]
Finally, Proposition 1.5 joined with Theorem 2.14 permits one to express $G_{q,j}(N,x)$ as a polynomial with integral coefficients in the $c^i(N)$, the $c^i(x)$, and $\varepsilon(x)$. One deduces the following theorem.

\textbf{Theorem 2.17.} Suppose the field $k$ has characteristic $0$, and let $i:Y\to X$ be an injection morphism of nonsingular quasi-projective varieties. Then one has formulas (i) to (v), and moreover, for every integer $j\geq 0$, the formula
\begin{equation}\tag{2.37}
 c^{q+j}(i_*(x))=
 i_*\left(\left(\frac{\widetilde c(x)*\lambda_{-1}(\widetilde c(N))}{(-1)^qN^q}\right)^{(j)}\right).
\end{equation}
The second member of formula (2.37) is an abbreviated way of denoting the polynomial in the components of $\widetilde c(x)$ and $\widetilde c(N)$ described in Proposition 1.5.

\textbf{Corollary.} Under the preceding conditions one has
\begin{equation}\tag{2.38}
 \ch(\widetilde c(i_*(x)))=i_*\left(\ch(\widetilde c(x))\,\mathfrak C(\check N)^{-1}\right),
\end{equation}
where $\ch$ is the Chern homomorphism defined by formula (1.26).

It is enough to use (2.37) written in the form
\begin{equation}\tag{2.37 bis}
 \widetilde c(i_*(x))=1+i_*\left(\frac{\widetilde c(x)*\lambda_{-1}(\widetilde c(N))-1}{(-1)^qN^q}\right),
\end{equation}
the explicit formula (1.29), formula (iv bis), and finally (1.30).

One may write (2.38) in the form
\begin{equation}\tag{2.38 bis}
 \ch_X(F)=i_*\left(\ch_Y(F)\,\mathfrak C(T_{X/Y})^{-1}\right),
\end{equation}
where $F$ is a coherent algebraic sheaf on $Y$ (considered in the first member as a sheaf on $X$ which is zero outside $Y$), where $\ch_X$ and $\ch_Y$ are the Chern homomorphisms taken respectively on $X$ and $Y$, and where $T_{X/Y}$ is the normal bundle of $Y$ in $X$.

\textbf{Remarks.} Contrary to (2.38), formula (2.37) is a ``formula without denominators''. It has been proved, even in characteristic $0$, only by placing oneself in $\Gr(K(X))$ and not in $A(X)$, which may not be the same thing. It is certainly true in every characteristic and in $A(X)$. Let us note that one verifies at once that the images of the two members by $i^*$ are the same; hence, applying $i_*$, the product of their difference by $N^q$ is zero. This suggests a completely different method of proof, by trying to obtain $Y$ as a transverse inverse image of a subvariety $Y'$ in a variety $X'$ such that, in dimension $\leq n=\dim X$, the class of $Y'$ in $A(X')$ is not a divisor of $0$. Unfortunately one cannot hope for $X'$ to be a Grassmannian or something close to it.

Note that, if torsion elements are neglected, formula (2.38) is equivalent to (2.37).

\end{document}
